
\documentclass[a4paper,12pt]{article}
\usepackage[margin=2cm]{geometry}

% Pacchetti utili di base
\usepackage[utf8]{inputenc}   % codifica dei caratteri
\usepackage[T1]{fontenc}      % font encoding
\usepackage[italian]{babel}   % lingua italiana
\usepackage{graphicx}         % immagini
\usepackage{hyperref}         % link cliccabili
\usepackage{amssymb}
\usepackage{float}
\usepackage{amsmath}
\usepackage{enumitem}

% Informazioni sul documento
\title{Basi di Dati}
\author{Nicolò Giuliani}
\date{}


\begin{document}

\maketitle
\tableofcontents
\newpage

\section{Introduzione}

\subsection{Sistema informativo e sistema informatico}

\textbf{Base di dati} (significato generico): insieme organizzato di dati che
supporta lo svolgimento delle attività specifiche di un ente (istituzione,
azienda, ufficio, persona).\\
\textbf{Base di dati} (significato specifico): insieme di dati gestito da un
DBMS.

\begin{description}[style=nextline]
  \item[Sistema informativo] componente di un ente che gestisce le informazioni
    di interesse. Ogni ente ha il proprio sistema informativo, anche se non
    esplicitamente identificabile nella sua organizzazione; va studiato nel suo
    contesto operativo. È un concetto \emph{indipendente} dall'automazione:
    esistono da secoli organizzazioni la cui unica missione è gestire dati
    (anagrafi, banche).
  \item[Sistema informatico] la porzione automatizzata del sistema informativo,
    cioè la componente che gestisce le informazioni con tecnologie
    informatiche.
\end{description}

Gerarchia di inclusione: Azienda $\supset$ Organizzazione $\supset$ Sistema
informativo $\supset$ Sistema informatico.

Gestione dell'informazione: raccolta/acquisizione, archiviazione/conservazione,
elaborazione/trasformazione/produzione, distribuzione/comunicazione/scambio.

\subsection{Informazione e dati}

Le attività umane rappresentano l'informazione in molti modi (idee informali,
linguaggio naturale, disegni e diagrammi, numeri e codici) su supporti diversi
(mente umana, carta, dispositivi elettronici). Nei sistemi informatici invece
l'informazione è rappresentata in modo povero ma essenziale: come \textbf{dati}.

\begin{itemize}
  \item \emph{informazione}: fatti forniti o appresi su qualcosa o qualcuno.
  \item \emph{dato}: elemento di informazione, punto di partenza fissato di un
    procedimento.
\end{itemize}

Il dato è il risultato del processo di organizzazione, codifica e memorizzazione
dell'informazione. \textbf{Il dato da solo è inutile senza la sua
interpretazione} (es.\ i numeri su un cartello stradale non dicono nulla se non
si sa che sono orari, e quale orario vale in quale giorno).

\paragraph{Perché i dati?} È difficile rappresentare con precisione informazione
e conoscenza ricche. I dati sono una \emph{risorsa strategica} perché sono più
stabili delle altre componenti (processi aziendali, tecnologie, ruoli): es.\ i
dati di banche e anagrafi.

\subsection{DBMS}

\textbf{DBMS} (DataBase Management System): sistema che gestisce collezioni di
dati
\begin{itemize}
  \item \textbf{grandi}
  \item \textbf{persistenti}
  \item \textbf{condivise}
\end{itemize}
garantendo
\begin{itemize}
  \item \textbf{privatezza}
  \item \textbf{affidabilità}
  \item \textbf{efficienza}
  \item \textbf{efficacia}
\end{itemize}

Esempi: DB2, Oracle, SQLServer, MySQL, PostgreSQL, Access, BigQuery, SQLite.

\paragraph{Grandi} dimensioni molto maggiori della memoria centrale; il solo
limite considerato è quello fisico dei dispositivi. Ordini di grandezza: 1--5 TB
(dati transazionali), 30--50 TB (dati decisionali), 500--800 TB (dati
scientifici), $10^{11}$ record.

\paragraph{Persistenti} il loro tempo di vita è indipendente da quello dei
processi che li usano.

\paragraph{Condivise} ogni (grande) azienda è divisa in settori, ognuno con un
proprio sotto-sistema informativo. Senza condivisione nascono
\textbf{ridondanza} (stessi dati ripetuti) e quindi possibile
\textbf{incoerenza} (descrizioni diverse non coincidono). La base di dati è una
risorsa \emph{integrata e condivisa}, da cui la necessità di
\textbf{meccanismi di autorizzazione} (attività diverse sugli stessi dati) e
\textbf{controllo della concorrenza} (più utenti sugli stessi dati).

\paragraph{Privatezza} si definiscono autorizzazioni del tipo ``l'utente A può
leggere tutti i dati e modificare il dato X''.

\paragraph{Affidabilità} resistenza a malfunzionamenti hardware e software. La
base di dati è una risorsa preziosa e va conservata a lungo termine. Tecnica
fondamentale: la gestione delle \textbf{transazioni}.

\paragraph{Efficienza} uso proficuo delle risorse di memoria (primaria e
secondaria) e di tempo (esecuzione e risposta). Poiché i DBMS offrono moltissime
funzioni, potrebbero risultare inefficienti: da qui grandi investimenti e
concorrenza fra prodotti. L'efficienza è anche frutto della qualità del
software.

\paragraph{Efficacia} rendere produttive le attività degli utenti, offrendo
funzionalità potenti e flessibili.

\subsection{Transazioni}

Una \textbf{transazione} è un insieme di operazioni non separabili
(\emph{atomiche}), corrette anche in presenza di concorrenza, con effetti
definitivi.

\begin{description}[style=nextline]
  \item[Atomicità] o si esegue tutta la sequenza o non se ne esegue nulla. Es.:
    ``sposta denaro dal conto A al conto B: o si preleva da A e si versa su B, o
    non si fa nulla''.
  \item[Concorrenza] transazioni concorrenti devono restare coerenti: due assegni
    sullo stesso conto incassati contemporaneamente non devono perdersi; due
    agenzie non devono prenotare due volte lo stesso posto sul treno.
  \item[Durabilità] il \texttt{commit} garantisce che il risultato finale sia
    registrato e conservato, anche in caso di concorrenza o guasti.
\end{description}

Due punti di vista sulle transazioni:
\begin{itemize}
  \item \emph{utente}: un programma che realizza una funzione utile, frequente e
    predefinita, con poche eccezioni previste (versamento in conto, emissione di
    un certificato di nascita, iscrizione all'anagrafe, prenotazione di un volo).
    Sono scritte in un linguaggio ospite, noto o \emph{ad hoc}.
  \item \emph{sistema}: sequenza indivisibile di operazioni (cfr.\ affidabilità).
\end{itemize}

\subsection{Evoluzione storica}

\begin{description}[style=nextline]
  \item[Anni '70 --- nessun DBMS] le applicazioni accedono direttamente ai file
    del sistema operativo; i dati comuni vengono duplicati nelle applicazioni.
  \item[Anni '80 --- primi DBMS] fra le applicazioni (Pascal, C) e i file si
    interpone il DBMS, che offre tabelle.
  \item[Anni '90 --- DBMS attivi] il comportamento procedurale comune viene
    condiviso fra le applicazioni. Nascono le \textbf{stored procedure}, che però
    non sono standardizzate e soffrono di \emph{impedance mismatch} con il
    linguaggio usato per esprimerle; si introducono quindi i \textbf{trigger},
    regole che modellano il comportamento procedurale condiviso e sono gestite
    dal DBMS stesso.
  \item[Anni 2000] architettura client/server su Internet: front-end
    (Javascript) sul client, Web Application Server e back-end (Java), database
    attivo con tabelle e trigger sul server.
  \item[Anni 2010] app mobili: front-end in Java, Objective-C, HTML5; il resto
    dell'architettura resta invariato.
\end{description}

\subsection{DBMS e file system}

Anche un file system può gestire grandi collezioni di dati, ma offre un accesso
a grana grossa: ``tutto o niente''. I DBMS estendono le funzionalità del file
system, offrendo più servizi in modo integrato.

\paragraph{Descrizione dei dati} ogni programma che non usa un DBMS contiene al
proprio interno la descrizione della struttura dei file che tratta: programmi
diversi possono avere descrizioni diverse, generando incoerenza fra le
rappresentazioni. Nei DBMS una porzione della base di dati (il \textbf{catalogo}
o \emph{dizionario}) contiene una descrizione centralizzata dei dati, condivisa
da tutte le applicazioni.

\subsection{Modelli dei dati}

Nei DBMS i dati sono descritti a diversi \textbf{livelli di astrazione}: i
programmi usano una rappresentazione di alto livello indipendente da quella di
livello più basso, così una modifica al livello basso non richiede di modificare
i programmi. Questo è ottenuto attraverso il concetto di \textbf{modello dei
dati}.

\textbf{Modello dei dati}: insieme di costrutti utilizzati per organizzare i dati
e descriverne il comportamento. Componente fondamentale: dare struttura ai dati
tramite \textbf{costruttori di tipo}. Es.: il \textbf{modello relazionale}
fornisce il costruttore \emph{relazione}, che permette di definire insiemi di
record omogenei.

\paragraph{Schema e istanza} in ogni base di dati esistono
\begin{itemize}
  \item lo \textbf{schema}: invariante nel tempo, descrive la struttura dei dati
    (aspetto \emph{intensionale}) --- es.\ l'intestazione delle tabelle;
  \item l'\textbf{istanza}: i valori effettivi, che possono cambiare rapidamente
    (aspetto \emph{estensionale}) --- es.\ il corpo delle tabelle.
\end{itemize}

\paragraph{Due tipologie di modelli}
\begin{description}[style=nextline]
  \item[Modelli logici] usati dai DBMS per organizzare e memorizzare i dati;
    usati dai programmi, indipendenti dalla rappresentazione fisica. Es.:
    relazionale, reticolare, gerarchico, a oggetti, XML.
  \item[Modelli concettuali] rappresentano i dati indipendentemente da ogni
    sistema specifico; descrivono concetti del mondo reale e si usano nella fase
    preliminare di progettazione. Il più diffuso è \textbf{Entity-Relationship}.
\end{description}

\subsection{Architettura di un DBMS}

Architettura semplificata (a due livelli): utente $\to$ \emph{schema logico}
$\to$ \emph{schema interno} $\to$ DB.

\begin{description}[style=nextline]
  \item[Schema logico] descrive la struttura della base di dati tramite il
    modello logico dei dati (es.\ la struttura delle tabelle).
  \item[Schema interno (o fisico)] rappresenta lo schema logico a livello di
    memorizzazione, con strutture dati \emph{raw} (file, record, puntatori).
\end{description}

Poiché la rappresentazione logica è indipendente da quella fisica --- una tabella
si usa ``così com'è'', indipendentemente dalla rappresentazione sottostante, che
può anche variare nel tempo --- nel corso si considera solo il livello logico.

\paragraph{Architettura standard ANSI/SPARC a tre livelli} aggiunge sopra lo
schema logico uno o più \textbf{schemi esterni}, ciascuno usato da un gruppo di
utenti:
\begin{itemize}
  \item \textbf{schema logico}: struttura della base di dati secondo il modello
    logico ``principale'';
  \item \textbf{schema interno}: rappresentazione a livello di memorizzazione;
  \item \textbf{schema esterno}: descrive una parte della base di dati in un
    modello logico (\emph{viste} parziali, tabelle derivate, anche da modelli
    diversi).
\end{itemize}

\paragraph{Indipendenza dei dati} è conseguenza della stratificazione;
l'accesso ai dati avviene sempre a livello esterno (che può coincidere con
quello logico).
\begin{description}[style=nextline]
  \item[Indipendenza fisica] livello logico ed esterno sono indipendenti dal
    livello fisico: una relazione è usata allo stesso modo qualunque sia la sua
    realizzazione fisica, che può cambiare senza modificare i programmi.
  \item[Indipendenza logica] il livello esterno è indipendente da quello logico:
    si può operare su una vista senza modificare il livello logico, e modifiche
    ``trasparenti'' allo schema logico lasciano inalterato lo schema esterno.
\end{description}

\subsection{Linguaggi per basi di dati}

Un altro aspetto dell'efficacia è la disponibilità di linguaggi e interfacce
diverse:
\begin{itemize}
  \item linguaggi testuali \emph{interattivi} (\textbf{SQL});
  \item comandi SQL immersi in un linguaggio \textbf{ospite} (Pascal, Java, C);
  \item comandi SQL immersi in un linguaggio \emph{ad hoc} con altre
    funzionalità (grafici, stampa di tabelle) --- es.\ Oracle PL/SQL;
  \item interfacce amichevoli, senza linguaggio testuale.
\end{itemize}

\paragraph{Esempio SQL interattivo} date le tabelle \texttt{Corsi(Corso,
Docente, Aula)} e \texttt{Aule(Nome, Edificio, Piano)}, per ottenere i corsi che
si tengono al piano terra:
\begin{verbatim}
SELECT Corso, Aula, Piano
FROM Aule, Corsi
WHERE Nome = Aula
  AND Piano = 'Terra'
\end{verbatim}

\paragraph{Separazione fra dati e programmi}
\begin{description}[style=nextline]
  \item[DML] \emph{data manipulation language}: interrogazione e aggiornamento
    delle \emph{istanze} della base di dati.
  \item[DDL] \emph{data definition language}: definizione degli \emph{schemi}
    (logico, esterno, fisico) e altre operazioni. Es.:
\end{description}
\begin{verbatim}
CREATE TABLE orario (
  corso    CHAR(20),
  docente  CHAR(20),
  aula     CHAR(4),
  ora      CHAR(5) )
\end{verbatim}

\subsection{Attori e utenti}

\begin{itemize}
  \item progettisti e realizzatori di DBMS;
  \item progettisti della base di dati e amministratori (\textbf{DBA});
  \item progettisti e programmatori di applicazioni;
  \item utenti:
    \begin{itemize}
      \item \emph{utenti finali} (terminalisti): eseguono operazioni predefinite
        (transazioni);
      \item \emph{utenti casuali}: eseguono operazioni non previste a priori,
        usando linguaggi interattivi.
    \end{itemize}
\end{itemize}

\paragraph{DBA} persona o gruppo responsabile del controllo e della gestione
centralizzata della base di dati (efficienza, affidabilità, autorizzazioni).
Spesso il DBA si occupa anche della progettazione, tranne che nei progetti
complessi.

\subsection{Pregi e difetti dei DBMS}

\paragraph{Pregi}
\begin{itemize}
  \item dati come risorsa comune, base di dati come modello della realtà;
  \item gestione dei dati centralizzata, standardizzabile ed economia di scala;
  \item disponibilità di servizi integrati;
  \item riduzione di ridondanze e incoerenze;
  \item indipendenza dei dati (favorisce sviluppo e manutenzione delle
    applicazioni).
\end{itemize}

\paragraph{Difetti}
\begin{itemize}
  \item costo dei prodotti e della loro adozione;
  \item funzionalità non separabili (con conseguente riduzione di efficienza).
\end{itemize}

\section{Il modello relazionale}

\subsection{Modelli logici}

Esistono tre modelli logici \emph{tradizionali}:
\begin{itemize}
  \item \textbf{gerarchico};
  \item \textbf{reticolare} (network);
  \item \textbf{relazionale}.
\end{itemize}
Più recenti sono il modello a \emph{oggetti} (poco diffuso) e i modelli su
\emph{XML} (considerati complementari al relazionale).

\paragraph{Caratteristica distintiva}
\begin{description}[style=nextline]
  \item[Gerarchico e reticolare] usano \textbf{riferimenti espliciti}
    (puntatori) fra i record.
  \item[Relazionale] è \textbf{value based}: i riferimenti fra dati memorizzati
    in relazioni diverse sono rappresentati tramite \emph{valori} (non
    puntatori).
\end{description}

\subsection{Il modello relazionale}

Definito da \textbf{E.\ F.\ Codd} nel 1970 per garantire l'indipendenza dei
dati. Implementato in DBMS reali solo a partire dal 1981: realizzare
indipendenza dei dati in modo efficiente e affidabile non è banale. Il modello
si basa sulla definizione logica di \emph{relazione} (con alcune differenze) e
le relazioni sono rappresentate tramite \textbf{tabelle}.

\paragraph{Relazione, relation, relationship} tre concetti da tenere distinti:
\begin{itemize}
  \item \emph{relazione logica}: nel senso della teoria degli insiemi;
  \item \emph{relazione} nel modello relazionale;
  \item \emph{relationship}: esprime una specifica classe di fatti nel modello
    Entity-Relationship.
\end{itemize}

\subsection{Relazione logica}

Dati $n$ insiemi $D_1, \ldots, D_n$ (non necessariamente distinti):
\begin{itemize}
  \item il \textbf{prodotto cartesiano} $D_1 \times \cdots \times D_n$ è
    l'insieme di tutte le n-uple $(d_1, \ldots, d_n)$ tali che $d_1 \in D_1,
    \ldots, d_n \in D_n$;
  \item una \textbf{relazione logica} su $D_1, \ldots, D_n$ è un
    \emph{sottoinsieme} $r \subseteq D_1 \times \cdots \times D_n$;
  \item i $D_i$ sono i \textbf{domini} della relazione.
\end{itemize}

\paragraph{Proprietà}
\begin{itemize}
  \item è un \emph{insieme} di n-uple: non c'è ordine fra le tuple;
  \item le tuple sono tutte \emph{distinte};
  \item ogni n-upla è \emph{ordinata}: l'$i$-esimo valore proviene
    dall'$i$-esimo dominio.
\end{itemize}

\paragraph{Struttura posizionale} in \texttt{Matches} $\subseteq$ \emph{string}
$\times$ \emph{string} $\times$ \emph{int} $\times$ \emph{int} lo stesso
dominio compare con \textbf{ruoli} diversi, distinti solo dalla \emph{posizione}
(1$^\text{a}$ colonna = squadra di casa, 2$^\text{a}$ = ospite, ecc.). Questa è
una struttura \textbf{posizionale}.

\subsection{Struttura non posizionale}

Nel modello relazionale a ciascun dominio si associa un \textbf{attributo}, un
nome unico che ne esprime il ``ruolo''. La posizione delle colonne diventa
irrilevante: la struttura è \textbf{non posizionale}. Es.\
\texttt{(Home, Away, GoalsH, GoalsA)}: scambiando l'ordine delle colonne (con
le relative intestazioni) l'informazione non cambia; scambiare \emph{solo} le
intestazioni cambia invece il significato.

\paragraph{Tabelle e relazioni} una tabella rappresenta una relazione se
\begin{itemize}
  \item tutte le righe sono \emph{diverse};
  \item tutte le intestazioni di colonna sono \emph{diverse};
  \item i valori all'interno di ogni colonna sono \emph{omogenei} (stesso
    dominio).
\end{itemize}
Righe e colonne possono assumere qualunque posizione.

\subsection{Modello value based}

I riferimenti fra dati in relazioni diverse sono espressi tramite \textbf{valori}
nelle tuple (es.\ \texttt{Exams.Student} = \texttt{Students.Number}), non tramite
puntatori. Altri modelli (reticolare, gerarchico, a oggetti) usano invece
riferimenti espliciti gestiti dal sistema.

\paragraph{Pregi del modello value based}
\begin{itemize}
  \item indipendenza dalla struttura fisica, che può cambiare dinamicamente
    (una cosa analoga si può ottenere con puntatori ``di alto livello'');
  \item si memorizza solo ciò che è rilevante dal punto di vista applicativo;
  \item utente e programmatore vedono gli stessi dati;
  \item i dati possono essere facilmente condivisi fra ambienti diversi;
  \item i puntatori sono direzionali, i valori no.
\end{itemize}

\subsection{Schemi e istanze}

\paragraph{Schema di una relazione} un nome di relazione $R$ con un insieme di
attributi $A_1, \ldots, A_n$:
\[
  R(A_1, \ldots, A_n)
\]
Es.\ \texttt{STUDENTS(Number, Surname, Name, Year of Birth)}.

\paragraph{Schema di database relazionale} insieme di schemi di relazione:
\[
  \mathcal{R} = \{R_1(X_1), \ldots, R_k(X_k)\}
\]

\paragraph{Tupla} una tupla su un insieme di attributi $X$ è un
\textbf{mapping} che associa a ogni attributo $A \in X$ un valore nel dominio
di $A$. Con $t[A]$ si indica il valore della tupla $t$ sull'attributo $A$
(es.\ $t[\texttt{Name}] = \text{Mario}$).

\paragraph{Istanze}
\begin{itemize}
  \item una \textbf{istanza di relazione} $r$ è un insieme finito di tuple con
    schema $R(X)$;
  \item una \textbf{istanza di database relazionale} su
    $\mathcal{R} = \{R_1(X_1), \ldots, R_n(X_n)\}$ è un insieme
    $r = \{r_1, \ldots, r_n\}$ di istanze di relazione ($r_i$ istanza di $R_i$).
\end{itemize}

\subsection{Casi particolari}

\paragraph{Relazioni con un solo attributo} sono ammesse: es.\
\texttt{WorkerStudents(Number)}, che elenca soltanto le matricole degli
studenti lavoratori.

\paragraph{Strutture annidate} molti dati reali (es.\ scontrini, fatture) sono
naturalmente \emph{annidati}: un'intestazione con vari campi ripetuti (righe di
prodotto). Il modello relazionale è \emph{piatto}: le strutture annidate si
rappresentano con più relazioni collegate da valori. Es.\
\texttt{Receipts(Number, Date, Total)} e
\texttt{Courses(Receipt, Row, No, Description, Price)}. La rappresentazione
scelta dipende dagli obiettivi: righe ripetute ammesse? ordine delle righe
significativo? esistono più rappresentazioni possibili.

\subsection{Informazione parziale}

Il modello relazionale ha una struttura \emph{rigida}: l'informazione è
espressa in tuple che devono rispettare uno schema. Ma i dati reali possono non
combaciare con il formato atteso (es.\ un ``secondo nome'' non sempre
esistente).

\paragraph{Perché non usare valori speciali} scegliere un valore del dominio (0,
stringa vuota, ``99'', $\ldots$) per marcare l'assenza è pericoloso:
\begin{itemize}
  \item potrebbero non esserci valori ``inutilizzati'' disponibili;
  \item quei valori potrebbero in seguito diventare significativi;
  \item bisognerebbe comunque tracciare a livello di software questi valori
    ``non rappresentabili''.
\end{itemize}

\paragraph{Valore NULL} tecnica grezza ma efficace: introdurre un valore
speciale \textbf{NULL} che indica l'assenza di un valore atteso e non appartiene
al dominio. Formalmente, per ogni attributo $A$, $t[A]$ è o un valore in
$\text{dom}(A)$ oppure \texttt{NULL}. Servono vincoli specifici per stabilire
quando \texttt{NULL} non è ammesso.

\paragraph{Tipi di NULL} ci sono almeno tre casi diversi di valore mancante:
\begin{itemize}
  \item valore \emph{sconosciuto} (esiste ma non lo conosciamo);
  \item valore \emph{inesistente} (non esiste per quella tupla);
  \item valore \emph{non informativo} (non sappiamo se esiste).
\end{itemize}
I DBMS moderni non distinguono fra questi casi.

\paragraph{Troppi NULL} una istanza con troppi \texttt{NULL} (specialmente nelle
chiavi o nei riferimenti) perde valore informativo: l'uso di \texttt{NULL} va
limitato.

\subsection{Vincoli di integrità}

Alcune istanze sintatticamente corrette possono non rappresentare informazione
corretta per l'applicazione (es.\ un voto ``32'', una \emph{laude} con voto
$<30$, un riferimento a uno studente inesistente).

\paragraph{Vincolo di integrità} proprietà che deve valere sulle istanze per
rappresentare informazione corretta. Formalmente è una funzione booleana
(\textbf{predicato}) che associa a ogni istanza \emph{true} o \emph{false}.

\paragraph{Perché}
\begin{itemize}
  \item descrizione più accurata del dominio reale;
  \item supporto alla \emph{data quality};
  \item utili in fase di progettazione;
  \item usati nella valutazione delle interrogazioni.
\end{itemize}

\paragraph{Osservazione} i DBMS non supportano tutti i tipi di vincolo: si
dichiarano quelli supportati e il DBMS ne impedisce la violazione; per gli
altri il programmatore deve implementarli fuori dal DBMS.

\paragraph{Classificazione}
\begin{description}[style=nextline]
  \item[Intra-relazionali] coinvolgono una sola relazione:
    \begin{itemize}
      \item su valori (\emph{domain constraints});
      \item su tuple.
    \end{itemize}
  \item[Inter-relazionali] coinvolgono più relazioni.
\end{description}

\subsection{Vincoli di tupla}

Esprimono regole sui valori di ogni singola tupla, indipendentemente dalle
altre. Caso specifico: i \textbf{domain constraints} coinvolgono un solo
attributo.

\paragraph{Sintassi} espressione booleana di atomi che confrontano:
\begin{itemize}
  \item valori del dominio degli attributi, es.\
    \texttt{(Grade $\geq$ 18) AND (Grade $\leq$ 30)};
  \item espressioni aritmetiche su tali valori, es.\
    \texttt{GrossPay = (Deductions + Net)}.
\end{itemize}
Una tupla \texttt{(Bruni, 50.000, 11.000, 36.000)} viola il vincolo perché
$11.000 + 36.000 \neq 50.000$.

\subsection{Chiavi}

Nel mondo reale ogni entità è identificabile: nella tabella \texttt{Students}
ogni tupla ha un \texttt{Number} unico e non esistono tuple con stessi
\texttt{Surname, Name, Date of Birth}. Formalizziamo:

\begin{description}[style=nextline]
  \item[Superchiave] un insieme $K$ di attributi di $r$ tale che non esistono
    due tuple distinte $t_1, t_2 \in r$ con $t_1[K] = t_2[K]$.
  \item[Chiave] una superchiave \emph{minimale}: nessun suo sottoinsieme
    proprio è già superchiave.
\end{description}

\paragraph{Esempi} in \texttt{Students(Number, Surname, Name, Curriculum,
Birth)}:
\begin{itemize}
  \item \texttt{Number} è chiave: è superchiave (identifica ciascuna tupla) e
    contiene un solo attributo, quindi è minimale;
  \item \texttt{\{Surname, Name, Birth\}} è un'altra chiave: è superchiave ed
    è minimale;
  \item \texttt{\{Surname, Curriculum\}} potrebbe essere superchiave in
    un'istanza particolare ma non nel caso generale.
\end{itemize}

\paragraph{Vincoli, schema e stati}
\begin{itemize}
  \item i vincoli corrispondono a proprietà del mondo reale, modellate nel DB;
  \item sono modellati sullo \emph{schema} (si applicano a tutte le tuple);
  \item ogni schema può avere un insieme di vincoli: le tuple memorizzate sono
    \textbf{corrette} (valide) se soddisfano tutti i vincoli;
  \item un'istanza può soddisfare \emph{per caso} altri vincoli non dichiarati.
\end{itemize}

\paragraph{Esistenza delle chiavi} ogni relazione:
\begin{itemize}
  \item non può contenere tuple con tutti i valori uguali (è un insieme);
  \item ha come superchiave l'insieme di tutti i suoi attributi;
  \item quindi ha almeno una chiave.
\end{itemize}

\paragraph{Importanza}
\begin{itemize}
  \item l'esistenza di almeno una chiave garantisce l'\emph{accessibilità} di
    ogni tupla;
  \item le chiavi permettono di correlare tuple fra relazioni diverse: il
    modello relazionale è \emph{value based}.
\end{itemize}

\subsection{Chiavi e valori nulli}

Se una tupla contiene \texttt{NULL} sugli attributi della chiave, i suoi
valori non permettono di:
\begin{itemize}
  \item identificare la tupla;
  \item creare facilmente riferimenti esterni da altre relazioni.
\end{itemize}
L'uso di \texttt{NULL} sulle chiavi va limitato.

\paragraph{Chiave primaria} fra le chiavi di una relazione se ne sceglie una
(\textbf{primary key}) che \emph{non ammette valori NULL}. Convenzionalmente si
indica con la sottolineatura degli attributi nello schema.

\subsection{Integrità referenziale}

Le informazioni fra relazioni diverse sono correlate tramite valori comuni, in
particolare valori di chiavi (primarie). Tali correlazioni devono essere
\emph{coerenti}: es.\ il codice del vigile in \texttt{Infringement.Policeman}
deve corrispondere a un \texttt{ID} esistente in \texttt{Policemen}.

\paragraph{Vincolo di integrità referenziale} (\emph{foreign key}) fra gli
attributi $X$ di una relazione $R_1$ e una relazione $R_2$: i valori di $X$
in $R_1$ devono comparire come valori della chiave primaria di $R_2$.

\paragraph{Esempio} in un DB con \texttt{Infringement(Code, Date, Policeman,
Country, Plate)}, \texttt{Policemen(ID, Surname, Name)},
\texttt{Car(Country, Plate, Surname, Name)}:
\begin{itemize}
  \item \texttt{Policeman} $\to$ \texttt{Policemen(ID)};
  \item \texttt{(Country, Plate)} in \texttt{Infringement} $\to$
    \texttt{Car(Country, Plate)}.
\end{itemize}
Una tupla \texttt{Infringement} con \texttt{(Italy, AJ046EZ)}, mentre in
\texttt{Car} \texttt{AJ046EZ} appartiene alla \texttt{France}, viola il
vincolo.

\paragraph{Osservazioni}
\begin{itemize}
  \item ruolo cruciale nel modello value based;
  \item i vincoli possono essere \emph{rilassati} in presenza di \texttt{NULL};
  \item le violazioni causate da operazioni possono essere gestite tramite
    \emph{azioni compensative};
  \item attenzione ai vincoli definiti su più attributi: l'\emph{ordine} degli
    attributi è rilevante (in \texttt{Crashes(Code, CountryA, PlateA, CountryB,
    PlateB)} servono due vincoli distinti verso \texttt{Car(Country, Plate)},
    uno per $(A,A)$ e uno per $(B,B)$).
\end{itemize}

\subsection{Azioni compensative}

Quando una modifica (es.\ cancellazione di una tupla referenziata) violerebbe
un vincolo referenziale, il DBMS può reagire in modi diversi.

\paragraph{Comportamento standard}
\begin{description}[style=nextline]
  \item[Restrict] rifiuta la cancellazione.
\end{description}

\paragraph{Azioni compensative possibili}
\begin{description}[style=nextline]
  \item[Cascade] rimuove ``a catena'' le tuple referenzianti (es.\ eliminando
    il progetto \texttt{XYZ} da \texttt{Projects} si eliminano anche i
    dipendenti in \texttt{Employees} con \texttt{Project = XYZ}).
  \item[Set NULL] pone \texttt{NULL} nell'attributo referenziante
    (\texttt{Employees.Project} diventa \texttt{NULL} per chi lavorava a
    \texttt{XYZ}).
  \item[Set default] introduce un valore di default.
\end{description}

\end{document}
